\documentclass[a4paper,11pt]{article}

\usepackage{fontspec}
\usepackage[ngerman,provide=*]{babel}
\usepackage{csquotes}

\usepackage[a4paper,margin=2.5cm,headheight=14pt]{geometry}
\usepackage{setspace}
\onehalfspacing
\usepackage{parskip}
\usepackage{booktabs}
\usepackage{tabularx}
\usepackage{array}
\usepackage{amsmath}
\usepackage{enumitem}
\setlist[itemize]{leftmargin=*,topsep=2pt,itemsep=2pt}
\setlist[enumerate]{leftmargin=*,topsep=2pt,itemsep=2pt}

\usepackage{fancyhdr}
\pagestyle{fancy}
\fancyhf{}
\renewcommand{\headrulewidth}{0.3pt}
\renewcommand{\footrulewidth}{0pt}
\fancyhead[L]{\footnotesize MTA -- Sitzung 2: Wie funktionieren Topic-Modelle?}
\fancyhead[R]{\footnotesize Seite \thepage}
\fancyfoot[C]{\footnotesize Prof.\ Dr.\ Christian Papilloud · Institut für Soziologie · MLU Halle-Wittenberg}

\usepackage{titlesec}
\titleformat{\section}{\large\bfseries}{\thesection.}{0.6em}{}
\titleformat{\subsection}{\normalsize\bfseries}{\thesubsection}{0.6em}{}
\titlespacing*{\section}{0pt}{1.5em}{0.5em}
\titlespacing*{\subsection}{0pt}{1.0em}{0.3em}

\usepackage{xcolor}
\definecolor{boxgrau}{gray}{0.95}
\usepackage[most]{tcolorbox}
\newtcolorbox{merkbox}[1][]{
  colback=boxgrau,colframe=black!40,boxrule=0.4pt,
  arc=2pt,left=8pt,right=8pt,top=4pt,bottom=4pt,
  fonttitle=\bfseries,#1
}

\usepackage{hyperref}
\hypersetup{
  colorlinks=true,
  linkcolor=black!70,
  urlcolor=black!60,
  citecolor=black!60,
  pdftitle={MTA Sitzung 2 -- Handout},
  pdfauthor={Christian Papilloud},
}

\title{\textbf{Wie funktionieren Topic-Modelle?}\\[4pt]
\large Sitzung 2 -- Mathematische Grundlagen, LDA, NMF\\
und der Sprung zu BERT/LLMs\\[8pt]
\normalsize Begleitendes Handout}
\author{Prof.\ Dr.\ Christian Papilloud\\
\small Institut für Soziologie -- Martin-Luther-Universität Halle-Wittenberg}
\date{Master Soziologie -- 2026}

\begin{document}

\maketitle
\thispagestyle{fancy}

\vspace{-1em}

% ===============================================================
\section{Wo wir stehen}

In der ersten Sitzung haben wir Topic-Modelle (TM) methodologisch
positioniert: \emph{quantitativ basiert}, aber im Dienst der
qualitativen Forschung, näher an Mayrings qualitativer Inhaltsanalyse
als an Oevermanns objektiver Hermeneutik. Heute fragen wir: \emph{Wie
funktioniert das technisch?} -- vom kleinen Beispiel bis zu den
neuesten Verfahren.

Dieses Handout folgt der gleichen Struktur wie die Sitzung:

\begin{enumerate}
\item Was ist ein \emph{topic}? Die Matrixfaktorisierung anhand eines
kleinen Beispiels.
\item Vergleich mit Clusteranalyse und Faktorenanalyse.
\item Die zwei klassischen Verfahren: LDA und NMF.
\item Die neue Welt: BERTopic und Large Language Models.
\end{enumerate}

% ===============================================================
\section{Was ist ein \emph{topic}?}

\subsection{Die Intuition}

Ein \emph{topic} ist nichts Geheimnisvolles. Es ist eine
\textbf{Wortliste} -- eine Gruppe von Wörtern, die in den
Dokumenten der untersuchten Textquelle \emph{oft gemeinsam}
vorkommen. Wenn die Wörter \enquote{Milch} und \enquote{Katze}
immer wieder in denselben Dokumenten auftauchen, gibt es vermutlich
ein topic, das beide enthält. Wenn ein drittes Wort, \enquote{Schale},
ebenfalls oft gemeinsam mit den ersten beiden auftaucht, gehört es
ebenfalls dazu. So entsteht durch wiederholte Ko-Okkurrenz eine
Wortgruppe, die wir als \emph{semantisches Thema} interpretieren.

Ein Topic-Modell liefert immer zwei Resultate:

\begin{itemize}
\item Die \textbf{Topic-Wort-Zugehörigkeit}: für jedes topic eine
Liste von Wörtern, jeweils mit einem \emph{Gewicht}, das ausdrückt,
wie zentral das Wort für dieses topic ist.

\item Die \textbf{Topic-Gewichte der Dokumente}: für jedes Dokument
eine Verteilung über die topics, die ausdrückt, in welchem Maße
das Dokument zu welchem topic gehört.
\end{itemize}

In MTA werden wir diese beiden Resultate als CSV-Tabellen, JSON-Dateien
und als Heatmap-Diagramme wiederfinden.

\subsection{Das MLU-Sommer-Beispiel}

Das Buch von Papilloud \& Hinneburg (2018, Abb.~2.1) illustriert die
Idee mit drei kleinen Dokumenten:

\begin{itemize}
\item \emph{Dokument 1:} Die Studenten der MLU können die Teile des
Campus im Sommer per Fahrrad erreichen.
\item \emph{Dokument 2:} Im Sommer gehen Menschen gern schwimmen oder
fahren Fahrrad.
\item \emph{Dokument 3:} Die MLU bietet auf ihrem Campus den Studenten
kostenlos Internet an.
\end{itemize}

Wir reduzieren die Texte auf ihre Substantive: \emph{Campus, Fahrrad,
Internet, Mensch, Sommer, Student, Teil, MLU}. Intuitiv erkennen Sie
zwei Themenfelder: \textbf{Universität} (MLU, Campus, Student, Internet)
und \textbf{Sommer} (Sommer, Fahrrad, Mensch).

Die zentrale Frage ist: Kann ein Algorithmus diese zwei Themen
\emph{ohne semantisches Vorwissen} finden? Die Antwort ist ja -- und
zwar durch eine mathematische Operation, die \textbf{Matrixfaktorisierung}.

\subsection{Schritt 1: die Dokument-Wort-Matrix}

Wir bauen eine Tabelle: jede Zeile ist ein Dokument, jede Spalte ein
Wort. Der Eintrag in Zeile $i$, Spalte $j$ ist die Häufigkeit, mit
der Wort $j$ in Dokument $i$ vorkommt:

\smallskip
\begin{center}
\footnotesize
\begin{tabular}{@{}l|cccccccc@{}}
\toprule
& Campus & Fahrrad & Internet & Mensch & Sommer & Student & Teil & MLU \\
\midrule
Dok.\ 1 & 1 & 1 & 0 & 0 & 1 & 1 & 1 & 1 \\
Dok.\ 2 & 0 & 1 & 0 & 1 & 1 & 0 & 0 & 0 \\
Dok.\ 3 & 1 & 0 & 1 & 0 & 0 & 1 & 0 & 1 \\
\bottomrule
\end{tabular}
\end{center}
\smallskip

Diese \textbf{Dokument-Wort-Matrix} -- nennen wir sie $A$ -- ist der
Input jedes Topic-Modells. Bei einem realen Korpus von z.B.\ zehntausend
Texten mit einem Vokabular von fünftausend Wörtern enthält sie 50
Millionen Einträge, von denen die allermeisten Nullen sind (kein
einzelner Text enthält das gesamte Vokabular). Solche Matrizen heißen
\emph{sparse}.

\subsection{Schritt 2: die Faktorisierung}

Ein Topic-Modell zerlegt $A$ in das \emph{Produkt zweier kleinerer
Matrizen}:
\[
\underbrace{A}_{\text{Dokumente}\times\text{Wörter}}
\;\approx\;
\underbrace{W}_{\text{Dokumente}\times K}
\;\cdot\;
\underbrace{H}_{K\times\text{Wörter}}
\]
wobei $K$ die Anzahl der topics ist -- ein Parameter, den wir vorher
festlegen (im Beispiel: $K=2$).

Was bedeuten $W$ und $H$? Die Matrix $W$ sagt für jedes Dokument, wie
stark es zu jedem topic gehört (die \emph{Topic-Gewichte}). Die Matrix
$H$ sagt für jedes topic, wie stark jedes Wort zu diesem topic gehört
(der \emph{Topic-Inhalt}).

\begin{merkbox}
\textbf{Warum funktioniert das?} Die Faktorisierung erzwingt eine
\emph{Datenkompression}: Statt 24 Einträge in $A$ zu speichern,
brauchen wir nur $3\times 2 + 2\times 8 = 22$ Einträge in $W$ und $H$.
Bei großen Matrizen ist die Reduktion enorm. Diese Kompression
\emph{zwingt} das Modell, die \textbf{wesentlichen Strukturen} der
Daten zu finden -- jene Strukturen, die sich \enquote{lohnen}, weil sie
viele Dokumente und viele Wörter zugleich erklären.
\end{merkbox}

Das Ergebnis im Beispiel:

\smallskip
\textbf{Topic-Wort-Zugehörigkeit $H$:}
\begin{center}
\footnotesize
\begin{tabular}{@{}l|cccccccc@{}}
\toprule
& Campus & Fahrrad & Internet & Mensch & Sommer & Student & Teil & MLU \\
\midrule
Topic 1 (\emph{Uni}) & 1 & 0 & 1 & 0 & 0 & 1 & 1 & 1 \\
Topic 2 (\emph{Sommer}) & 0 & 1 & 0 & 1 & 1 & 0 & 0 & 0 \\
\bottomrule
\end{tabular}
\end{center}

\smallskip
\textbf{Topic-Gewichte $W$:} Dok.\ 1 enthält \emph{beide} topics
(Gewicht 1 und 1). Dok.\ 2 nur Sommer (0 und 1). Dok.\ 3 nur Uni
(1 und 0). Der Algorithmus hat ohne semantisches Vorwissen die zwei
Themenfelder gefunden, die wir intuitiv erkannt hatten.

\subsection{Wie findet der Algorithmus die Zerlegung?}

Das Problem ist \emph{zirkulär}: Wenn man die topics kennt, kann man
die Dokumente aufteilen; wenn man die Aufteilung kennt, kann man die
topics bestimmen. Aber zu Beginn kennt man keines von beiden. Die
Lösung ist die \textbf{iterative Annäherung}:

\begin{enumerate}
\item Man startet mit einer \emph{zufälligen} Konfiguration der topics.
\item Man berechnet daraus die beste Dokumentaufteilung.
\item Man berechnet daraus wieder die besten topics.
\item Man wiederholt, bis sich nichts mehr ändert (Konvergenz).
\end{enumerate}

Mit jeder Iteration verringert sich der \emph{Rekonstruktionsfehler}
-- der Abstand zwischen der originalen Matrix $A$ und ihrer
Rekonstruktion $W \cdot H$. Weil dieser Fehler nach unten beschränkt
ist (er kann nicht kleiner als 0 werden), konvergiert das Verfahren.

\subsection{Eine wichtige Eigenschaft: die \enquote{Halluzinationen}}

Wenn wir $W$ und $H$ multiplizieren, erhalten wir eine
\emph{rekonstruierte} Version der Dokument-Wort-Matrix. Diese
Rekonstruktion ist nicht identisch mit dem Original.

Im Beispiel hatte Dok.\ 1 ursprünglich keinen Eintrag für die Wörter
\enquote{Internet} und \enquote{Mensch}. Die Rekonstruktion sagt aber:
diese Wörter \emph{passen} zu Dok.\ 1, denn dieses Dokument enthält
die topics, zu denen sie gehören.

\begin{merkbox}
\textbf{Diese \enquote{Halluzinationen} sind kein Fehler.} Sie sind
einer der wichtigsten qualitativen Mehrwerte von Topic-Modellen:
Sie machen \emph{implizite} Bedeutungen sichtbar -- semantische
Nähen, die im Text \emph{vorhanden}, aber nicht \emph{explizit
ausgesprochen} sind. Wenn ein Text über \enquote{Studenten} und
\enquote{Campus} spricht, ohne \enquote{Internet} zu erwähnen, sagt
uns das Modell trotzdem, dass dieses Thema in der Nähe ist.
\end{merkbox}

% ===============================================================
\section{Vergleich mit Cluster- und Faktorenanalyse}

\subsection{Topic-Modelle vs.\ Clusteranalyse}

Auf den ersten Blick wirken Topic-Modelle wie eine Spielart der
Clusteranalyse. Beide bilden Gruppen von Objekten. Aber es gibt einen
fundamentalen Unterschied:

\begin{itemize}
\item Eine Clusteranalyse (etwa K-Means) teilt Dokumente
\textbf{disjunkt} ein: jedes Dokument gehört zu einem und nur einem
Cluster. Das ist eine \emph{harte} Klassifikation.

\item Topic-Modelle sind \textbf{Mixed-Membership-Models}
(Airoldi et al.\ 2014): Ein Dokument ist eine \emph{Mischung} aus
topics. Es kann zu 60\,\% aus Topic 1, zu 30\,\% aus Topic 2 und zu
10\,\% aus Topic 3 bestehen.
\end{itemize}

Dieser Unterschied ist nicht nur technisch. In der
sozialwissenschaftlichen Praxis ist er entscheidend: Texte sind in
aller Regel \emph{thematisch heterogen}. Ein Interview enthält
meistens mehrere Themen; ein wissenschaftlicher Artikel ebenso; eine
Zeitungsausgabe sowieso. Eine harte Clusterung würde dieser
Heterogenität nicht gerecht.

\textbf{Ein anschauliches Beispiel:} Drei Gruppen von Dokumenten --
nur über Sommer (Gruppe A), nur über die Universität (Gruppe B), über
beides (Gruppe C). Eine Clusteranalyse fände \emph{drei} Cluster:
Sommer, Uni, Sommer+Uni. Die Kombination wird zu einem eigenen
Cluster. Ein Topic-Modell findet \emph{zwei} topics: Sommer und Uni.
Die Gruppe C wird durch eine Mischung der beiden topics beschrieben.

Topic-Modelle finden also weniger, aber \emph{interpretierbarere}
Strukturen. Sie \emph{zerlegen} die Komplexität in elementarere
Themen, statt sie durch Kombinationen zu vervielfachen. Mathematisch
ausgedrückt (Erosheva, Fienberg \& Joutard 2007, später Galyardt 2014):
Ein Topic-Modell mit $K$ topics, das Dokumente mit insgesamt $T$
Wörtern beschreibt, kann durch ein Cluster-Modell mit $K^T$ Clustern
\emph{repräsentiert} werden. Die Mixed-Membership-Repräsentation ist
also \emph{viel} kompakter.

\subsection{Topic-Modelle vs.\ Faktorenanalyse}

Auch die Faktorenanalyse (FA) und die Hauptkomponentenanalyse (PCA)
sind Matrixfaktorisierungen: $X \approx Z \cdot W^T$. Sie sehen
strukturell wie Topic-Modelle aus. Es gibt aber zwei Unterschiede:

\textbf{Erstens: die Vorzeichen.} PCA und FA lassen \emph{negative}
Einträge in $Z$ und $W$ zu. Ein negativer Wert bedeutet semantisch:
\enquote{dieses Wort soll \emph{nicht} im Dokument vorkommen}. Solche
\enquote{negativen topics} sind für die qualitative Textanalyse kaum
interpretierbar. Topic-Modelle wie NMF und LDA verlangen, dass
\emph{alle Einträge} $\geq 0$ sind. Diese Einschränkung macht die
gefundenen topics zu \emph{additiven Beiträgen}, die sich klar und
eindeutig interpretieren lassen.

\textbf{Zweitens: die Verteilungsannahmen.} PCA setzt implizit eine
\emph{Normalverteilung} der Daten voraus. Worthäufigkeiten sind aber
nicht normalverteilt: sie folgen Long-Tail-Verteilungen mit sehr
vielen Nullen und einigen extrem häufigen Wörtern. Topic-Modelle wie
LDA nutzen passendere Verteilungen (Multinomial, Dirichlet-Multinomial),
die diese Eigenschaft besser abbilden.

% ===============================================================
\section{Die zwei klassischen Verfahren: LDA und NMF}

\subsection{LDA -- Latent Dirichlet Allocation}

LDA (Blei, Ng \& Jordan 2003) ist ein \emph{probabilistisches}
Topic-Modell. Es beschreibt die Erzeugung von Dokumenten als
Zufallsprozess. Die Geschichte, die das Modell erzählt, klingt so:

\begin{enumerate}
\item Jedes topic ist eine \emph{Wahrscheinlichkeitsverteilung über
das Vokabular}. Wenn man dem topic eine Münze in die Hand gibt und
sie würfeln lässt, kommt mit hoher Wahrscheinlichkeit ein Wort heraus,
das zentral zum topic gehört -- mit niedriger Wahrscheinlichkeit ein
peripheres Wort. Parameter: $\varphi_k$ für topic $k$.

\item Jedes Dokument hat ebenfalls eine \emph{Verteilung über die
topics}. Wenn man dem Dokument eine Münze in die Hand gibt, kommt
mit hoher Wahrscheinlichkeit dasjenige topic, das dieses Dokument
dominiert. Parameter: $\theta_d$ für Dokument $d$.

\item \emph{Wie entsteht ein Wort im Dokument?} Für jedes Wort:
würfle zuerst ein topic aus $\theta_d$, dann würfle aus diesem topic
ein Wort aus seiner Verteilung $\varphi_k$.
\end{enumerate}

Das ist natürlich nicht, wie Menschen tatsächlich Texte schreiben.
Es ist ein \emph{generatives Modell}, das mathematisch gut zu
handhaben ist. Die Berechnung kehrt das Modell um: gegeben die Wörter
im Korpus, was sind die wahrscheinlichsten topics $\varphi$ und
Dokument-Verteilungen $\theta$?

Mathematisch sucht man die \emph{Posterior-Verteilung}
$p(\varphi, \theta, z \mid w)$. Sie ist nicht direkt berechenbar; es
gibt zwei Approximationsverfahren: die \textbf{Variationsinferenz}
(Blei et al.\ 2003) und das \textbf{Collapsed Gibbs Sampling}
(Griffiths \& Steyvers 2004). Was Sie sich merken sollten:
\emph{LDA liefert Wahrscheinlichkeiten, keine Punktwerte.}

\subsection{NMF -- Non-Negative Matrix Factorization}

NMF (Lee \& Seung 1999, 2001) löst dasselbe Problem ohne probabilistische
Mathematik, rein durch lineare Algebra. Die Aufgabe ist ein
Optimierungsproblem:
\[
\min_{W, H \geq 0} \| A - W \cdot H \|_F
\]
wobei $\| \cdot \|_F$ die \emph{Frobenius-Norm} ist -- eine Maßzahl,
die misst, wie weit die rekonstruierte Matrix $W \cdot H$ vom
Original $A$ entfernt ist. Der Algorithmus, der diese Optimierung
löst, heißt \emph{Alternating Least Squares}: man hält $W$ fest und
optimiert $H$, dann hält man $H$ fest und optimiert $W$, und so weiter,
bis Konvergenz erreicht ist.

Im Gegensatz zu LDA hat NMF \emph{keine} probabilistische Interpretation.
Aber NMF ist meistens \emph{schneller}, \emph{deterministischer}
(bei festem Startwert: identische Ergebnisse) und in vielen praktischen
Situationen \emph{stabiler}, besonders bei kleinen Korpora.

\subsection{Wann LDA, wann NMF?}

\begin{tabularx}{\textwidth}{@{}p{0.22\textwidth}XX@{}}
\toprule
& \textbf{LDA} & \textbf{NMF} \\
\midrule
Grundlage
  & Probabilistisches Modell
  & Lineare Algebra \\
\addlinespace[3pt]
Ausgabe
  & Wahrscheinlichkeiten
  & Gewichte (positiv) \\
\addlinespace[3pt]
Reproduzierbarkeit
  & Stochastisch, variiert je Lauf
  & Deterministisch bei festem Startwert \\
\addlinespace[3pt]
Geschwindigkeit
  & langsamer
  & schneller \\
\addlinespace[3pt]
Kleine Korpora
  & instabil
  & oft besser \\
\addlinespace[3pt]
Theoretische Eleganz
  & Bayessches Netzwerk
  & Optimierungsproblem \\
\bottomrule
\end{tabularx}

In MTA sind beide Verfahren implementiert. Eine empfehlenswerte
Praxis ist, \emph{beide laufen zu lassen} und ihre Ergebnisse zu
vergleichen. Wenn LDA und NMF ähnliche topics liefern, ist das ein
starkes Indiz für die Stabilität der gefundenen Strukturen. Wenn sie
sehr unterschiedliche topics liefern, ist Vorsicht geboten -- entweder
die Topic-Anzahl ist falsch gewählt, oder das Korpus ist zu klein, oder
die Vorverarbeitung muss verbessert werden.

\subsection{Wie viele topics?}

Sowohl LDA als auch NMF setzen die Anzahl $K$ der topics als
\emph{externen Parameter} voraus. Diese Wahl ist nicht trivial; sie
ist methodologisch und interpretativ.

Drei Strategien sind in der Literatur und in MTA verfügbar:

\begin{enumerate}
\item \textbf{Perplexity / Log-Likelihood} (für LDA): Wie gut beschreibt
das Modell \emph{neue} Daten, die nicht für das Training verwendet
wurden? Theoretisch elegant, aber: nicht jedes topic, das die
Perplexity verbessert, ist \emph{interpretierbar} (Chang et al.\ 2009).
Menschliche Bewertungen korrelieren manchmal sogar \emph{negativ}
mit Perplexity-Werten.

\item \textbf{Topic-Kohärenz} (PMI, Newman et al.\ 2010): Wie oft
tauchen die top-Wörter eines topics nahe beieinander im Korpus auf?
Diese Maßzahl korreliert besser mit menschlicher Bewertung. Verschiedene
Varianten existieren (Röder, Both \& Hinneburg 2015).

\item \textbf{Kreuzvalidierung mit Clusterverfahren} (in MTA
implementiert): MTA berechnet eine Reihe von Indikatoren -- Elbow,
Silhouette, Calinski-Harabasz, Davies-Bouldin, Cophenet -- und schlägt
Ihnen einen Wertebereich für $K$ vor.
\end{enumerate}

\begin{merkbox}
\textbf{Keine dieser Methoden liefert \emph{die eine} richtige Zahl.}
Sie liefern einen \emph{Wertebereich}, in dem die Entscheidung
interpretativ getroffen werden muss. Die Frage \enquote{wie viele
topics?} ist immer auch eine \emph{theoretische} Frage: Wie fein
soll Ihre Themenstruktur sein? Was wollen Sie sehen?
\end{merkbox}

% ===============================================================
\section{Die neue Welt: BERTopic und LLMs}

\subsection{Ein methodologischer Sprung seit 2018}

Das Buch von Papilloud \& Hinneburg ist 2018 erschienen. Damals
galten LDA und NMF als \emph{state of the art}. Seitdem hat sich das
Feld erheblich verändert. Vier Schritte:

\begin{itemize}
\item \textbf{2013 -- Word2Vec} (Mikolov et al.): die Idee der
\emph{word embeddings}. Jedes Wort wird durch einen Vektor in einem
hochdimensionalen Raum repräsentiert (typischerweise 100 bis 300
Dimensionen). Wörter mit ähnlicher Bedeutung haben \emph{ähnliche}
Vektoren. Die berühmte Gleichung
$\vec{v}(\text{König}) - \vec{v}(\text{Mann}) + \vec{v}(\text{Frau})
\approx \vec{v}(\text{Königin})$
illustriert das Prinzip.

\item \textbf{2018-2019 -- BERT} (Devlin et al.): kontextuelle
Embeddings. Dasselbe Wort hat verschiedene Vektoren je nach Satzkontext.
\enquote{Bank} in \enquote{Sitzbank im Park} ist nicht
\enquote{Bank} in \enquote{Geld auf der Bank}.

\item \textbf{2022 -- BERTopic} (Grootendorst 2022): ein
Topic-Modell, das BERT-Embeddings nutzt -- der direkte
\enquote{Nachfolger} von LDA und NMF im Stand der Forschung.

\item \textbf{2022 ff.\ -- ChatGPT, GPT-4, Claude, LLaMA \dots}: Large
Language Models, die Texte direkt \emph{verstehen} und
\emph{annotieren} können, ohne dass man ihnen eine
Dokument-Wort-Matrix vorlegt.
\end{itemize}

\subsection{Word2Vec, BERT und das Konzept der Embeddings}

\emph{Embeddings} sind die zentrale Neuerung, die alles, was danach
kommt, ermöglicht. Die Grundidee: anstatt Wörter als \emph{Symbole}
zu behandeln (\enquote{Hund} ist einfach das fünfte Wort im Vokabular,
mit Index 5), behandelt man sie als \emph{Punkte in einem geometrischen
Raum}. In diesem Raum gilt: semantisch ähnliche Wörter sind räumlich
nahe.

Klassische Topic-Modelle (LDA, NMF) arbeiten mit \emph{rohen}
Worthäufigkeiten. Für sie sind die Wörter \enquote{König} und
\enquote{Königin} unverbunden -- sie sehen sich nur die Ko-Okkurrenz
in Dokumenten an. Embeddings lösen dieses Problem: \enquote{König}
und \enquote{Königin} haben ähnliche Vektoren, weil sie in ähnlichen
Kontexten auftreten.

In MTA finden Sie diese Idee bereits in der Funktion \emph{Semantic
Context}: dort werden Word2Vec-ähnliche Embeddings (in der Standardform
über Ko-Okkurrenz und PCA) genutzt, um semantische Nachbarn von Wörtern
zu finden. Die Idee der Embeddings ist also keine ferne
Zukunftsmusik -- sie ist bereits Teil unserer Praxis.

\subsection{BERTopic -- ein moderner Topic-Modell-Workflow}

BERTopic (Grootendorst 2022) bricht mit der Matrixfaktorisierung. Sein
Workflow besteht aus vier Schritten:

\begin{enumerate}
\item Jedes Dokument wird durch \textbf{BERT} in einen
hochdimensionalen Vektor verwandelt (typischerweise 768 Dimensionen).

\item Diese Vektoren werden durch \textbf{UMAP} auf wenige Dimensionen
reduziert (Vorbereitung für das Clustering).

\item Die reduzierten Vektoren werden durch \textbf{HDBSCAN}
geclustert. Jeder Cluster ist ein topic.

\item Für jeden Cluster wird die top-Wortliste durch eine angepasste
TF-IDF-Heuristik berechnet (\textbf{c-TF-IDF}, class-based TF-IDF).
\end{enumerate}

Statt eine Matrix zu faktorisieren, ordnet BERTopic also Dokumente
einem Cluster zu -- es ist eigentlich näher an einer Clusteranalyse
als an einem klassischen Topic-Modell. Das macht es schnell und
kontextsensitiv, hat aber Konsequenzen.

\subsection{Empirische Vergleiche -- die Forschungslage}

Mehrere empirische Vergleiche der vergangenen Jahre haben LDA, NMF
und BERTopic gegenübergestellt (Egger \& Yu 2022; Kaur \& Wallace
2024; Mishra 2026; Ma et al.\ 2025).

\textbf{Wo BERTopic besser abschneidet:}

\begin{itemize}
\item Bei \emph{kurzen} Texten (Tweets, SMS, Antworten auf offene
Fragen), wo Ko-Okkurrenz allein nicht ausreicht.
\item In semantisch \emph{nuancierten} Domänen, wo verschiedene
Bedeutungen desselben Wortes wichtig sind.
\item Höhere \emph{Topic-Kohärenz} in vielen Benchmarks; auch
qualitative Forschende präferieren oft BERTopic-Ergebnisse
(Kaur \& Wallace 2024 zeigen 67\,\% Präferenz für BERTopic in
einer Studie mit 12 Forschenden).
\end{itemize}

\textbf{Wo LDA und NMF kompetitiv bleiben:}

\begin{itemize}
\item Bei \emph{längeren} Texten (Artikeln, Interviews, Büchern),
wo die Ko-Okkurrenz im Dokument reichlich Information liefert.
\item Wenn \emph{Interpretierbarkeit} und \emph{Reproduzierbarkeit}
wichtiger sind als marginale Kohärenz-Gewinne.
\item Bei \emph{kleinen} Korpora -- BERTopic verlangt oft sehr viele
Dokumente und verwirft regelmäßig 30 bis 70\,\% als
\enquote{Outliers} (Soledad et al.\ 2022).
\item Bei \emph{spezialisierten} Sprachen oder Disziplinen, für die
kein guter vortrainierter BERT-Modell verfügbar ist (etwa
historisches Deutsch, juristische Texte einer bestimmten
Tradition, Korpora in selteneren Sprachen).
\end{itemize}

\subsection{Large Language Models als analytische Werkzeuge}

Seit 2022 verändert sich das Feld noch einmal. LLMs (GPT-4, Claude,
LLaMA \dots) bieten Möglichkeiten, die klassische TM nicht hatten:

\begin{itemize}
\item \textbf{Direkte Codierung von Texten} -- inklusive deduktiver
Codierung anhand eines Codebuchs (Tai et al.\ 2024).
\item \textbf{Automatische Beschriftung} bestehender topics
(TopicGPT, QualIT bei Amazon).
\item \textbf{Reflexive Begleitung}: das LLM schreibt Memos, weist
auf Widersprüche zwischen Codierungen hin, schlägt theoretische
Rahmungen vor (Sinha et al.\ 2024).
\item \textbf{Thematische Analyse von Interviews} ohne klassisches
TM (De Paoli 2023).
\end{itemize}

Eine umfangreiche Literatur ist seit 2023 entstanden. Hayes (2025) und
Than et al.\ (2025) bieten gute Überblicke aus soziologischer Sicht.

\subsection{Eine empirische Warnung: Ashwin, Chhabra \& Rao (2025)}

Eine besonders wichtige Arbeit ist 2025 in den \emph{Sociological
Methods \& Research} erschienen. Die Autoren untersuchten Interviews
mit Rohingya-Geflüchteten und ihren bengalischen Gastgebern in
Bangladesch. Sie verglichen die Codierung durch einen LLM mit der
Codierung durch trainierte menschliche Codiererinnen.

\begin{merkbox}
\textbf{Der zentrale Befund.} Die LLM-Codierungen sind nicht nur
\emph{ungenau}, sondern \textbf{systematisch verzerrt}. Die Fehler
des LLMs verteilen sich \emph{nicht zufällig}, sondern korrelieren
mit Merkmalen der Befragten -- ihrer Sprache, ihrem Bildungsgrad,
ihrer Herkunftsregion. Inferenzen, die auf solchen Codierungen
beruhen, können in eine \emph{bestimmte Richtung} systematisch falsch
sein.
\end{merkbox}

Die Konsequenz, die die Autoren ziehen: Es sei oft besser, ein
\emph{eigenes, kleineres} Modell auf hochwertigen, manuell codierten
Daten zu trainieren als einen generischen LLM einzusetzen. Diese
Empfehlung ist auch für die qualitative Sozialforschung relevant:
Quantitative Effizienz darf nicht auf Kosten methodologischer
Belastbarkeit gehen.

\subsection{Warum bleibt MTA bei LDA und NMF?}

Eine berechtigte Frage. Die Antwort hat vier Aspekte:

\begin{enumerate}
\item \textbf{Lokal, offline, autonom.} MTA läuft auf Ihrem Laptop
ohne Internetverbindung und ohne API-Schlüssel. Ihre Forschungsdaten
verlassen Ihre Maschine nicht. Das ist nicht nur eine Frage der
Bequemlichkeit, sondern auch der \emph{Forschungsethik} und des
\emph{Datenschutzes} -- besonders bei sensiblen Daten (Interviews
mit Geflüchteten, medizinische Texte, persönliche Korrespondenz).

\item \textbf{Interpretierbarkeit.} LDA und NMF liefern
nachvollziehbare Wortlisten und Gewichte. Sie können einer Kollegin
oder einer Gutachterin erklären, \emph{warum} ein topic so aussieht,
wie es aussieht. Bei BERTopic und LLMs ist die Kette zwischen Input
und Output deutlich opaker.

\item \textbf{Reproduzierbarkeit.} NMF mit festem Startwert produziert
\emph{exakt dieselben} Ergebnisse, jedes Mal. LDA ist stochastisch,
aber durch wiederholte Läufe lässt sich die Variabilität abschätzen.
LLMs sind dagegen oft schlecht reproduzierbar -- ihre Antworten
hängen vom Modellzustand, von der \emph{Temperature}, von kleinen
Variationen im Prompt ab.

\item \textbf{Forschungsdidaktik.} Wer LDA und NMF verstanden hat,
versteht den \emph{Kern} der Topic-Modellierung -- die
Matrixfaktorisierung, das Mixed-Membership-Prinzip, die Logik der
Ko-Okkurrenz. Wer mit BERTopic oder einem LLM anfängt, ohne
Embeddings zu verstehen, arbeitet mit einer Blackbox. Diese
Vorlesungen wollen Ihnen die Werkzeuge in die Hand geben, später
\emph{informiert} zwischen den Verfahren wählen zu können.
\end{enumerate}

% ===============================================================
\section{Fragen zum Mitnehmen}

Drei Fragen, über die Sie bis zur nächsten Sitzung nachdenken können:

\begin{enumerate}
\item Wir haben gesehen, dass Topic-Modelle \enquote{Halluzinationen}
produzieren: sie schreiben einem Dokument Wörter zu, die es
ursprünglich nicht enthielt. Wir haben diese Eigenschaft als
\emph{qualitativen Mehrwert} dargestellt. Unter welchen Bedingungen
könnte sie umgekehrt zum \emph{Problem} werden?

\item Ashwin et al.\ (2025) zeigen, dass LLMs systematisch verzerrte
Codierungen produzieren. Können Sie sich Quellen \emph{ähnlicher}
Verzerrungen vorstellen, die auch bei LDA oder NMF auftreten könnten?

\item Die Wahl der Topic-Anzahl $K$ ist immer auch eine
\emph{theoretische} Wahl. Wenn Sie Interviews mit Studierenden zum
Thema \enquote{digitale Lehre} analysieren würden, welche Theorien
würden Ihre Wahl von $K$ leiten?
\end{enumerate}

% ===============================================================
\section{Ausblick auf Sitzung 3}

In der dritten und letzten theoretischen Sitzung wenden wir uns dem
\emph{Vorher} und \emph{Nachher} des Algorithmus zu:

\begin{itemize}
\item \emph{Vor} dem Algorithmus: die Vorverarbeitung der Textquelle
(Bereinigung, Stoppwörter, TF-IDF, Lemmatisierung, Named Entity
Recognition). Was MTA \emph{tut} und was es \emph{nicht} tut.
\item Metadaten und Topic-Modelle: drei Strategien (Spaltung,
Hinzufügung, Lernergänzung).
\item \emph{Nach} dem Algorithmus: der kritische Blick (Buch, Kap.\ 8).
Was Topic-Modelle \emph{nicht} können. Zwei verbreitete
Erwartungen, die nicht erfüllt werden.
\item Was sich für unsere Forschungspraxis ändert -- und wie diese
drei theoretischen Sitzungen in die folgende Arbeit mit MTA münden.
\end{itemize}

% ===============================================================
\section*{Zitierte Literatur}
\addcontentsline{toc}{section}{Zitierte Literatur}

\footnotesize

\textsc{Grundlage:}

Papilloud, C., \& Hinneburg, A.\ (2018). \emph{Einführung in die
qualitative Analyse von Texten mit Topic-Modellen}. Wiesbaden:
Springer. \emph{Bes.\ Kap.\ 2 und 3.}

\bigskip
\textsc{Klassische Topic-Modelle und Matrixfaktorisierung:}

Airoldi, E.\ M., Blei, D.\ M., Erosheva, E.\ A., \& Fienberg, S.\ E.\
(Hrsg., 2014). \emph{Handbook of Mixed Membership Models and Their
Applications}. CRC Press.

Blei, D.\ M.\ (2012). Probabilistic Topic Models.
\emph{Communications of the ACM}, 55(4), 77--84.

Blei, D.\ M., Ng, A.\ Y., \& Jordan, M.\ I.\ (2003). Latent Dirichlet
Allocation. \emph{Journal of Machine Learning Research}, 3, 993--1022.

Chang, J., Boyd-Graber, J., Wang, C., Gerrish, S., \& Blei, D.\ M.\
(2009). Reading Tea Leaves: How Humans Interpret Topic Models.
\emph{NIPS}.

Erosheva, E.\ A., Fienberg, S.\ E., \& Joutard, C.\ (2007). Describing
Disability Through Individual-Level Mixture Models for Multivariate
Binary Data. \emph{Annals of Applied Statistics}, 1(2), 502--537.

Griffiths, T.\ L., \& Steyvers, M.\ (2004). Finding Scientific Topics.
\emph{PNAS}, 101(suppl.\ 1), 5228--5235.

Lee, D.\ D., \& Seung, H.\ S.\ (2001). Algorithms for Non-Negative
Matrix Factorization. In \emph{Advances in Neural Information
Processing Systems}.

Newman, D., Lau, J.\ H., Grieser, K., \& Baldwin, T.\ (2010).
Automatic Evaluation of Topic Coherence. \emph{NAACL HLT}.

Röder, M., Both, A., \& Hinneburg, A.\ (2015). Exploring the Space
of Topic Coherence Measures. \emph{WSDM}.

\bigskip
\textsc{Word Embeddings und BERT:}

Devlin, J., Chang, M.-W., Lee, K., \& Toutanova, K.\ (2019). BERT:
Pre-training of Deep Bidirectional Transformers for Language
Understanding. \emph{NAACL HLT}.

Mikolov, T., Sutskever, I., Chen, K., Corrado, G.\ S., \& Dean, J.\
(2013). Distributed Representations of Words and Phrases and their
Compositionality. \emph{NIPS}.

Mostafavi, M., Porter, M.\ D., \& Robinson, D.\ T.\ (2025). Contextual
Embeddings in Sociological Research: Expanding the Analysis of
Sentiment and Social Dynamics. \emph{Sociological Methodology}.

\bigskip
\textsc{BERTopic und Vergleiche:}

Egger, R., \& Yu, J.\ (2022). A Topic Modeling Comparison Between
LDA, NMF, Top2Vec, and BERTopic to Demystify Twitter Posts.
\emph{Frontiers in Sociology}, 7.

Grootendorst, M.\ (2022). BERTopic: Neural Topic Modeling with a
Class-Based TF-IDF Procedure. \emph{arXiv:2203.05794}.

Kaur, A., \& Wallace, J.\ R.\ (2024). Moving Beyond LDA: A Comparison
of Unsupervised Topic Modelling Techniques for Qualitative Data
Analysis of Online Communities. \emph{arXiv:2412.14486}.

Ma, L., et al.\ (2025). AI-powered Topic Modeling: Comparing LDA
and BERTopic in Analyzing Opioid-Related Cardiovascular Risks in
Women. \emph{Experimental Biology and Medicine}.

Mishra, M.\ (2026). Comparative Analysis of BERTopic Versus LDA:
Topic Modelling in Marketing Research. \emph{Asia-Pacific Journal of
Management Research and Innovation}.

\bigskip
\textsc{Large Language Models in der qualitativen Sozialforschung:}

Ashwin, J., Chhabra, A., \& Rao, V.\ (2025). Using Large Language
Models for Qualitative Analysis can Introduce Serious Bias.
\emph{Sociological Methods \& Research}.

Bail, C.\ A.\ (2024). Can Generative AI Improve Social Science?
\emph{PNAS}, 121(21).

De Paoli, S.\ (2023). Performing an Inductive Thematic Analysis of
Semi-Structured Interviews With a Large Language Model: An
Exploration and Provocation on the Limits of the Approach.
\emph{Social Science Computer Review}.

Hayes, A.\ S.\ (2025). \enquote{Conversing} With Qualitative Data:
Enhancing Qualitative Research Through Large Language Models (LLMs).
\emph{International Journal of Qualitative Methods}.

Sinha, R., et al.\ (2024). Generative AI as a Reflexive Co-Researcher
in Grounded Theory.

Tai, R.\ H., Bentley, L.\ R., Xia, X., et al.\ (2024). An Examination
of the Use of Large Language Models to Aid Analysis of Textual Data.
\emph{International Journal of Qualitative Methods}.

Than, N., Fan, L., Law, T., Nelson, L.\ K., \& McCall, L.\ (2025).
Updating \enquote{The Future of Coding}: Qualitative Coding with
Generative Large Language Models. \emph{Sociological Methods \&
Research}.

\end{document}
